\section*{Chapter \ref{ch:g12:reaction-rates} --- Reaction Rates}

\begin{solution}{exo:g12:reaction-rates:1}
Temperature; contact area; concentration; temperature.
\end{solution}

\begin{solution}{exo:g12:reaction-rates:2}
\qty{13.9}{min} (blue) and \qty{6.9}{min} (red). After two half-lives,
$10.0 / 4 = \qty{2.5}{mmol/L}$.
\end{solution}

\begin{solution}{exo:g12:reaction-rates:3}
Slope $(5.0 - 1.0)/10 = \qty{0.40}{mmol.L^{-1}.min^{-1}}$: that is the
\omterm{def:g12:reaction-rates:rate}{rate of appearance} at \qty{5}{min}.
\end{solution}

\begin{solution}{exo:g12:reaction-rates:4}
Cooling and diluting make the reaction so slow that the composition no
longer changes during the \omterm{def:g11:titration:titration}{titration}. The composition measured is the one
the \omterm{def:g6:pure-substances-and-mixtures:mixture}{mixture} had when the reaction was stopped, at the moment of pouring.
\end{solution}

\begin{solution}{exo:g12:reaction-rates:5}
\qty{30}{min} is 3 half-lives: $1/8$ remains. \qty{1}{h} is 6 half-lives:
$1/64$, about \qty{1.6}{\percent}.
\end{solution}

\begin{solution}{exo:g12:reaction-rates:6}
$\ln[\ce{A}]$: 2.079, 1.733, 1.386, 1.040, 0.693. It falls by 0.346 every
\qty{5}{min}: a straight line, first order. $k = 0.346/5 =
\qty{0.0693}{min^{-1}}$ and $t_{1/2} = \ln 2 / k = \qty{10.0}{min}$ (as
seen directly: 8.00, 4.00, 2.00 every \qty{10}{min}).
\end{solution}

\begin{solution}{exo:g12:reaction-rates:7}
$k = 0.693 / 25 = \qty{0.0277}{s^{-1}}$.
\end{solution}

\begin{solution}{exo:g12:reaction-rates:8}
$0.100\,e^{-0.20} = \qty{0.0819}{mol/L}$; $0.100\,e^{-0.60} =
\qty{0.0549}{mol/L}$; $0.100\,e^{-2.0} = \qty{0.0135}{mol/L}$.
\end{solution}

\begin{solution}{exo:g12:reaction-rates:9}
At \qty{13.9}{min}, $[\ce{A}] = \qty{5.0}{mmol/L}$; the tangent's slope is
close to $-0.25$, and $-k[\ce{A}] = -0.050 \times 5.0 =
\qty{-0.25}{mmol.L^{-1}.min^{-1}}$: half the initial value.
\end{solution}

\begin{solution}{exo:g12:reaction-rates:10}
Iodine is the only coloured species: by the Beer--Lambert law
$A = \varepsilon \ell [\ce{I2}]$. The \omterm{def:g12:reaction-rates:rate}{rate of appearance} is the slope of
the tangent to the curve $[\ce{I2}](t)$, that is the slope of the
tangent to $A(t)$ divided by $\varepsilon \ell$.
\end{solution}

\begin{solution}{exo:g12:reaction-rates:11}
The initial rate $k[\ce{A}]_0$ is twice as large for the second run; the
half-lives are equal, $\ln 2 / k$, independent of $[\ce{A}]_0$.
\end{solution}

\begin{solution}{exo:g12:reaction-rates:12}
$t = \ln(1/f)/k$. In half-lives: $\ln(1/f)/\ln 2$. For \qty{1}{\percent},
$\ln 100 / \ln 2 = 6.6$; for \qty{0.1}{\percent}, $\ln 1000 / \ln 2 =
10.0$.
\end{solution}

\begin{solution}{exo:g12:reaction-rates:13}
\qty{90}{\percent} used, $f = 0.10$: $t = \ln 10 / k$, that is
\qty{46}{min} at \qty{20}{\celsius} and \qty{23}{min} at \qty{30}{\celsius}.
The \omterm{def:g12:reaction-rates:first-order}{rate constant} doubles over \qty{10}{\celsius}: the rule holds here.
\end{solution}

\begin{solution}{exo:g12:reaction-rates:14}
The temperature drops close to \qty{0}{\celsius}, and the concentrations
are divided by 10. With a rate proportional to the product of two
concentrations, \omterm{def:g10:concentration-and-dilution:dilution}{dilution} alone divides it by $10 \times 10 = 100$.
\end{solution}

\begin{solution}{exo:g12:reaction-rates:15}
$v(t) = -\dfrac{d}{dt}\big([\ce{A}]_0 e^{-kt}\big) = k[\ce{A}]_0 e^{-kt}$.
Initial rate $k[\ce{A}]_0$; the rate is halved when $e^{-kt} = 1/2$, that
is after one \omterm{def:g12:reaction-rates:half-life}{half-life}.
\end{solution}

\begin{solution}{pb:g12:reaction-rates:1}
\textbf{1.} By the Beer--Lambert law, the dye being the only species
absorbing at this wavelength; the solution must be dilute enough
($A$ below about 1).

\textbf{2.} It absorbs no visible light: it adds nothing to $A$.

\textbf{3.} So that its concentration hardly changes: the rate then
depends on the dye alone.

\textbf{4.} $A = 0$ (the blank).

\textbf{5.} $A$ falls from 0.800 to 0.402 in \qty{6}{min}: about
\qty{6}{min}.

\textbf{6.} From 0.402 at \qty{6}{min} to 0.199 at \qty{12}{min}: halved
again in \qty{6}{min}.

\textbf{7.} $-0.223$, $-0.448$, $-0.691$, $-0.911$, $-1.155$, $-1.366$,
$-1.614$, $-1.945$, $-2.551$, $-3.079$.

\textbf{8.} The points lie on a straight line: the reaction is first
order in the dye.

\textbf{9.} $(-2.551 + 0.223)/20 = \qty{-0.116}{min^{-1}}$.

\textbf{10.} $k = \qty{0.116}{min^{-1}}$.

\textbf{11.} $0.693 / 0.116 = \qty{6.0}{min}$, as read in the table.

\textbf{12.} $k A_0 = 0.116 \times 0.800 = \qty{0.093}{min^{-1}}$ in
\omterm{def:g11:absorbance:absorbance}{absorbance} units per minute.

\textbf{13.} The same, \qty{6.0}{min}: a first-order \omterm{def:g12:reaction-rates:half-life}{half-life} does not
depend on the starting value.

\textbf{14.} $0.800\, e^{-0.116 \times 30} = 0.025$.

\textbf{15.} $t = \ln 100 / k = 4.61 / 0.116 = \qty{40}{min}$.

\textbf{16.} $40 / 6.0 = 6.6$ half-lives.

\textbf{17.} $t = \ln(0.800/0.010)/k = 4.38 / 0.116 = \qty{38}{min}$.

\textbf{18.} Still 6.6 half-lives, now $6.6 \times 3.0 = \qty{20}{min}$.

\textbf{19.} About 6.6 half-lives, that is about \qty{40}{min} here.
\end{solution}
